% 第5章第5节：赋范空间（范数与收敛、有限维赋范空间、完备性与Banach空间），只讨论有限维情形。

\section{赋范空间}
\label{sec:ls-normed-spaces}

前四节只用到加法与数乘，回答的都是“能否表示、是否唯一、有几个方向”这类是非与计数问题，还没有“大小”。可是用一个多项式近似一个函数时，要问误差有多大；一列近似越来越好时，要问它们是否趋于某个对象，那个对象是否仍在原来的空间里；上一节的读取也要问，输入不大时读数最多有多大。第~\ref{chap:linear-operations}章用范数度量$\mathbb R^d$中数组的大小（第~\ref{sec:ls-norms-distances}节），本节把它推广到一般线性空间。

本节回答三个问题：范数怎样给出收敛；有限维空间中，选哪个范数会不会影响收敛与否；收敛序列的极限什么时候一定留在空间中。前两个问题在有限维中有干净的答案：一切范数等价，收敛就是逐坐标收敛。第三个问题引出完备性与Banach空间，它在有限维中自动成立。本节只讨论有限维空间；无限维空间中这三个问题的答案都会改变，留到后面的章节。

\subsection{范数与赋范空间}
\label{sec:ls-normed-space-definition}

\paragraph{范数的定义}
第~\ref{chap:linear-operations}章的$\ell_1$、$\ell_2$、$\ell_\infty$范数各有侧重，却满足同样的几条要求：大小非负，只有零向量大小为零，整体放大$c$倍时大小放大$|c|$倍，两段误差相加时总大小不超过各自大小之和。这些要求只用到加法与数乘，因此可以直接作为一般线性空间上“大小”的要求。
\begin{definition}[赋范空间]
\label{def:ls-normed-space}
设$V$是实线性空间。映射$\lVert\cdot\rVert:V\to[0,\infty)$若对任意$\bm x,\bm y\in V$与$c\in\mathbb R$满足
\begin{enumerate}
 \item 正定性：$\lVert\bm x\rVert=0$当且仅当$\bm x=\bm0$；
 \item 齐次性：$\lVert c\bm x\rVert=|c|\,\lVert\bm x\rVert$；
 \item 三角不等式：$\lVert\bm x+\bm y\rVert\leq\lVert\bm x\rVert+\lVert\bm y\rVert$，
\end{enumerate}
就称为$V$上的范数；配备范数的线性空间$(V,\lVert\cdot\rVert)$称为\term{赋范空间}（normed space）。
\end{definition}
范数诱导距离$d(\bm x,\bm y)=\lVert\bm x-\bm y\rVert$：它非负，只在$\bm x=\bm y$时为零，对称，并满足$d(\bm x,\bm z)\leq d(\bm x,\bm y)+d(\bm y,\bm z)$。这些性质也是\BookChapterRef{chap:metric-spaces}{第~\ref{chap:metric-spaces}章}度量空间的公理，所以赋范空间是一类特殊的度量空间。它的距离还与线性运算相容：平移不改变距离，缩放把距离按比例缩放。

把三角不等式用于$\bm x=(\bm x-\bm y)+\bm y$，得$\lVert\bm x\rVert-\lVert\bm y\rVert\leq\lVert\bm x-\bm y\rVert$；交换$\bm x,\bm y$的角色，又得$\lVert\bm y\rVert-\lVert\bm x\rVert\leq\lVert\bm x-\bm y\rVert$。两式合起来就是反向三角不等式
\begin{equation}
 \bigl|\lVert\bm x\rVert-\lVert\bm y\rVert\bigr|\leq\lVert\bm x-\bm y\rVert,
 \label{eq:ls-reverse-triangle}
\end{equation}
它说明两个向量相距很近时，长度也相差很少。

\paragraph{例子}
$\mathbb R^d$上的$p$范数、同形矩阵上的Frobenius范数与诱导矩阵范数都是范数，第~\ref{chap:linear-operations}章已经验证。一般空间上，最直接的办法是借坐标：若$\lVert\cdot\rVert_\star$是$\mathbb R^d$上的范数，$\mathcal B$是$V$的基，则$\lVert\bm v\rVert=\lVert[\bm v]_{\mathcal B}\rVert_\star$是$V$上的范数，因为坐标映射是保持线性组合的双射。也可以不经过坐标，直接用对象本身的性质规定范数。下面在$P_k$上各举一例。
\begin{example}[多项式空间上的两个范数]
\label{ex:ls-polynomial-norms}
给定：$P_k$上的两个规则：对$p=a_0+a_1t+\cdots+a_kt^k$，
\[
 p\mapsto\max_{0\leq j\leq k}|a_j|,
 \qquad
 p\mapsto\lVert p\rVert_\infty=\max_{t\in[0,1]}|p(t)| .
\]

验证：第一个规则是坐标$(a_0,\ldots,a_k)$在基$(1,t,\ldots,t^k)$下的$\ell_\infty$范数，按上面的借坐标方法是范数。对第二个规则，齐次性逐点成立；三角不等式来自$|p(t)+q(t)|\leq|p(t)|+|q(t)|\leq\lVert p\rVert_\infty+\lVert q\rVert_\infty$对每个$t$成立，再对$t$取最大值；正定性来自次数不超过$k$的非零多项式至多有$k$个根，不能在整个区间$[0,1]$上为零。

结论：二者都是$P_k$上的范数，前者经过坐标，后者直接用函数值。
\end{example}

\paragraph{收敛}
有了距离，就可以说一列向量趋于某个向量。在赋范空间中，若实数列$\lVert\bm x_n-\bm x\rVert\to0$，就说$\bm x_n$按该范数\emph{收敛}到$\bm x$，记为$\bm x_n\to\bm x$。极限若存在则唯一：若$\bm x_n\to\bm x$且$\bm x_n\to\bm y$，则$\lVert\bm x-\bm y\rVert\leq\lVert\bm x-\bm x_n\rVert+\lVert\bm x_n-\bm y\rVert\to0$。由式~\eqref{eq:ls-reverse-triangle}，$\bm x_n\to\bm x$时还有$\lVert\bm x_n\rVert\to\lVert\bm x\rVert$。

收敛的定义依赖所选的范数：不同范数给出的误差读数不同，同一个序列会不会在一种范数下收敛、在另一种范数下不收敛？下一小节说明，有限维空间中不会出现这种情形。

\subsection{有限维赋范空间}
\label{sec:ls-finite-dimensional-sequences}
\label{sec:linear-algebra-sequence-convergence}
\label{sec:ls-error-scale-comparison}

讨论向量序列的收敛，先要选一个范数测量误差。$\ell_1$、$\ell_2$、$\ell_\infty$或某个加权范数给出的误差读数各不相同。本小节证明，在有限维空间中任意两个范数最多相差常数倍，因而收敛与否、极限是什么都与范数无关；证明要把第~\ref{chap:function-limits}章有界实数列的Bolzano--Weierstrass定理逐坐标用于向量序列。由此还得到线性映射的连续性、子空间的闭性与最近点的存在。

\paragraph{闭集与有界集}
在赋范空间中，若存在$R$使$\lVert\bm x\rVert\leq R$对集合$K$中的一切$\bm x$成立，就称$K$\emph{有界}；若$K$中任何收敛序列的极限仍属于$K$，就称$K$\emph{闭}。在实轴上，这一闭性与第~\ref{chap:function-limits}章用补集定义的闭集一致（引理~\ref{lem:analysis-closed-sequence-criterion}）；一般度量空间中的开闭集在\BookChapterRef{chap:metric-spaces}{第~\ref{chap:metric-spaces}章}讨论。

闭球$\{\bm x:\lVert\bm x\rVert\leq r\}$是闭集：若$\lVert\bm x_n\rVert\leq r$且$\bm x_n\to\bm x$，则$\lVert\bm x\rVert=\lim_n\lVert\bm x_n\rVert\leq r$。同理，单位球面$\{\bm x:\lVert\bm x\rVert=1\}$是闭有界集。若$K\subseteq V$上的函数$f$满足：只要$K$中的$\bm x_n\to\bm x\in K$，就有$f(\bm x_n)\to f(\bm x)$，就称$f$在$K$上连续。由式~\eqref{eq:ls-reverse-triangle}，范数本身是连续函数。

\paragraph{范数等价}
现在比较同一个有限维空间上的两个范数。
\begin{theorem}[有限维范数等价]
\label{thm:ls-norm-equivalence}
有限维实线性空间上任意两个范数$p,q$都等价：存在常数$c,C>0$，使$c\,q(\bm x)\leq p(\bm x)\leq C\,q(\bm x)$对所有$\bm x$成立。
\end{theorem}
\begin{proof}
思路是把每个范数都与一个固定的参照范数比较。固定基$(\bm b_1,\ldots,\bm b_d)$，把$\bm x=\sum_jx_j\bm b_j$等同于坐标向量$(x_1,\ldots,x_d)$，取坐标最大范数$\lVert\bm x\rVert_c=\max_j|x_j|$作参照。按它收敛就是逐坐标收敛：$\lVert\bm x_n-\bm x\rVert_c\to0$当且仅当每个坐标$x_{n,j}\to x_j$，因为最大值趋于零当且仅当有限个坐标差都趋于零。

上界。对任意范数$p$，由三角不等式与齐次性，
\[
 p(\bm x)=p\Bigl(\sum_jx_j\bm b_j\Bigr)\leq\sum_j|x_j|\,p(\bm b_j)\leq C_p\lVert\bm x\rVert_c,
 \qquad C_p=\sum_jp(\bm b_j).
\]

下界。由反向三角不等式，$|p(\bm x)-p(\bm y)|\leq p(\bm x-\bm y)\leq C_p\lVert\bm x-\bm y\rVert_c$，所以只要$\lVert\bm x_n-\bm x\rVert_c\to0$，就有$p(\bm x_n)\to p(\bm x)$。考虑坐标单位球面$S=\{\bm x:\lVert\bm x\rVert_c=1\}$，记$m_p=\inf_{\bm x\in S}p(\bm x)\geq0$，对每个$n$取$\bm x_n\in S$使$p(\bm x_n)<m_p+1/n$。这列向量的每个坐标都满足$|x_{n,j}|\leq1$，是有界实数列。由Bolzano--Weierstrass定理~\ref{thm:analysis-bounded-subsequence}，先在第一坐标上抽出收敛子列，再在这个子列中按第二坐标继续抽取，如此重复$d$次，最后得到的子列$\bm x_{n_k}$在每个坐标上都收敛。记各坐标的极限组成的向量为$\bm x^\star$，则$\lVert\bm x_{n_k}-\bm x^\star\rVert_c\to0$；由$\lVert\cdot\rVert_c$的反向三角不等式，$\lVert\bm x^\star\rVert_c=\lim_k\lVert\bm x_{n_k}\rVert_c=1$，即$\bm x^\star\in S$。再由上面$p$的连续性，\phantomsection\label{thm:ls-extreme-value}$p(\bm x^\star)=\lim_kp(\bm x_{n_k})=m_p$，即$p$在$S$上取到最小值$m_p$。由$\lVert\bm x^\star\rVert_c=1$知$\bm x^\star\neq\bm0$，范数的正定性保证$m_p=p(\bm x^\star)>0$。对任意$\bm x\neq\bm0$，$\bm x/\lVert\bm x\rVert_c$在$S$上，由齐次性得$p(\bm x)\geq m_p\lVert\bm x\rVert_c$。合起来，
\begin{equation}
 m_p\lVert\bm x\rVert_c\leq p(\bm x)\leq C_p\lVert\bm x\rVert_c.
 \label{eq:linear-algebra-finite-norm-equivalence}
\end{equation}

最后比较$p$与$q$：$p\leq C_p\lVert\cdot\rVert_c\leq(C_p/m_q)\,q$，交换$p,q$的角色得反向不等式。
\end{proof}
定理中的常数依赖维数与所选范数。看$\ell_1$与$\ell_\infty$这一对。
\begin{example}[$\ell_1$与$\ell_\infty$范数的等价常数]
\label{ex:ls-l1-linf-constants}
给定：$\mathbb R^d$上的$\lVert\bm x\rVert_1=\sum_j|x_j|$与$\lVert\bm x\rVert_\infty=\max_j|x_j|$。

验证：$d$个非负数$|x_1|,\ldots,|x_d|$之和不小于其中最大的一个，又不超过最大者的$d$倍，所以
\[
 \lVert\bm x\rVert_\infty\leq\lVert\bm x\rVert_1\leq d\lVert\bm x\rVert_\infty .
\]
取$\bm x=\bm e_1$，两个范数都为$1$，左边取等号；取$\bm x=(1,\ldots,1)^\top$，$\lVert\bm x\rVert_1=d$、$\lVert\bm x\rVert_\infty=1$，右边取等号。

结论：常数$1$与$d$都不能改进；维数越大，同一误差在两种范数下的读数可以相差越多。
\end{example}
证明中真正用到有限维的是下界：坐标只有有限个，逐坐标抽取$d$次才得到在每个坐标上都收敛的子列，从而在单位球面上取到正的最小值。

\paragraph{收敛与连续性}
由范数等价，$\lVert\bm x_n-\bm x\rVert\to0$当且仅当$\lVert\bm x_n-\bm x\rVert_c\to0$，即每个坐标$x_{n,j}\to x_j$。因此在有限维空间中，序列是否收敛、收敛到哪里，都与所选范数和基无关，后文可以按论证方便选用范数。矩阵序列看作$\mathbb R^{m\times n}$中的向量，结论相同：$\bm A_k\to\bm A$在任意矩阵范数下都等价于逐元素收敛。

闭集与有界集同样与范数无关：两个范数只差常数倍，一个范数下有界的集合在另一个范数下也有界，收敛序列也完全相同。于是下界证明中的逐坐标抽取对任何范数都适用：有界序列的坐标有界，可以抽出收敛子列；若序列取自闭集，极限仍在集合内。

由此得到闭有界集上连续函数的最值。设$f$是非空闭有界集$K$上的连续函数，$m=\inf_{K}f$。取$K$中的序列$\bm x_n$，使$f(\bm x_n)\to m$（若$m=-\infty$，则使$f(\bm x_n)\to-\infty$），从中抽出收敛子列$\bm x_{n_k}\to\bm x^\star$；由闭性，$\bm x^\star\in K$。由连续性，$f(\bm x_{n_k})\to f(\bm x^\star)$。若$m=-\infty$，这个有限极限与$f(\bm x_{n_k})\to-\infty$矛盾，所以$m$有限；于是$f(\bm x^\star)=\lim_kf(\bm x_{n_k})=m$，即下确界可以取到。上确界同理。下面几处“可以达到”都由这一论证得到。

线性映射也自动连续。设$T:V\to W$是有限维赋范空间之间的线性映射，固定$V$的基并沿用证明中的记号。由三角不等式与齐次性，
\[
 \lVert T\bm x\rVert_W\leq\sum_j|x_j|\,\lVert T\bm b_j\rVert_W\leq C\lVert\bm x\rVert_c\leq\frac C{m_V}\lVert\bm x\rVert_V,
 \qquad C=\sum_j\lVert T\bm b_j\rVert_W,
\]
其中$m_V>0$是式~\eqref{eq:linear-algebra-finite-norm-equivalence}中$\lVert\cdot\rVert_V$的下界常数。于是$\lVert T\bm x_n-T\bm x\rVert_W\leq(C/m_V)\lVert\bm x_n-\bm x\rVert_V$，$T$连续。

上式还说明比值$\lVert T\bm x\rVert_W/\lVert\bm x\rVert_V$有界。由上面的最值论证，连续函数$\bm x\mapsto\lVert T\bm x\rVert_W$在闭且有界的单位球面上取到最大值，所以第~\ref{chap:linear-operations}章诱导矩阵范数定义中的上确界有限且可以达到。取$W=\mathbb R$，就回答了上一节留下的问题：线性泛函$\varphi$的大小可以定义为
\[
 \lVert\varphi\rVert_*=\sup_{\lVert\bm x\rVert\leq1}|\varphi(\bm x)|,
\]
它有限并且取到，是对偶空间$V^*$上的范数；在$\mathbb R^d$上，$\varphi(\bm x)=\bm z^\top\bm x$时它正是第~\ref{sec:ls-dual-norms}节的对偶范数$\lVert\bm z\rVert_*$。

\paragraph{子空间的闭性与商范数}
\phantomsection
\label{sec:ls-quotient-norms}
点到子空间的距离为零，是否意味着点就在子空间里？这取决于子空间是否闭。
\begin{theorem}[有限维子空间闭性]
\label{thm:ls-subspace-closed}
有限维赋范空间的每个线性子空间$U$都是闭集：若$\bm x_n\in U$且$\bm x_n\to\bm x$，则$\bm x\in U$。
\end{theorem}
\begin{proof}
把$U$的基扩充为整个空间的基，则$\bm x_n$在补充方向上的坐标恒为零。收敛等价于逐坐标收敛，所以$\bm x$在补充方向上的坐标是零序列的极限，仍为零，于是$\bm x\in U$。
\end{proof}
记点到子空间的距离为$d(\bm x,U)=\inf_{\bm u\in U}\lVert\bm x-\bm u\rVert$。若$d(\bm x,U)=0$，可在$U$中取$\bm u_n$使$\lVert\bm x-\bm u_n\rVert\to0$，由闭性$\bm x\in U$；反之显然。所以$d(\bm x,U)=0$当且仅当$\bm x\in U$。

这一结论保证了商空间上也能度量大小。第~\ref{sec:ls-quotient-spaces}节的商空间$V/U$把相差$U$中向量的点看成同一个类$\bm x+U$，那里只建立了线性结构。同一类中各点到$U$的距离相同：若$\bm x'=\bm x+\bm u_0$，$\bm u_0\in U$，则$\bm u$遍历$U$时$\bm u-\bm u_0$也遍历$U$，所以$d(\bm x',U)=d(\bm x,U)$。因此可以把这个距离当作类的大小，称为商范数
\[
 \lVert\bm x+U\rVert_{V/U}=d(\bm x,U)=\inf_{\bm u\in U}\lVert\bm x-\bm u\rVert .
\]
它满足范数的三条要求。三角不等式：对任意$\bm u,\bm u'\in U$，$d(\bm x+\bm y,U)\leq\lVert(\bm x-\bm u)+(\bm y-\bm u')\rVert\leq\lVert\bm x-\bm u\rVert+\lVert\bm y-\bm u'\rVert$，再对$\bm u,\bm u'$分别取下确界。齐次性：$c\neq0$时$cU=U$，所以$d(c\bm x,U)=|c|\,d(\bm x,U)$；$c=0$时两边都为零。正定性要求$d(\bm x,U)=0$只在$\bm x+U=U$即$\bm x\in U$时发生，这正是上面的闭性结论。

\paragraph{一般范数下的最近点}
距离$d(\bm y,U)$定义为下确界，它总能取到吗？由$d(\bm y,U)\leq\lVert\bm y-\bm0\rVert=\lVert\bm y\rVert$，而$\lVert\bm u\rVert>2\lVert\bm y\rVert$时$\lVert\bm y-\bm u\rVert\geq\lVert\bm u\rVert-\lVert\bm y\rVert>\lVert\bm y\rVert$，所以只需在$K=\{\bm u\in U:\lVert\bm u\rVert\leq2\lVert\bm y\rVert\}$上找最小值。由子空间闭性，$K$是闭有界集，按“收敛与连续性”一段的最值论证，连续函数$\bm u\mapsto\lVert\bm y-\bm u\rVert$在其上取到最小值，所以最近点存在。

唯一性却可能失败。
\begin{example}[反例：$\ell_\infty$范数下最近点不唯一]
\label{ex:ls-linf-nonunique-nearest}
给定：$\mathbb R^2$取$\ell_\infty$范数，子空间$U=\operatorname{span}\{(1,0)^\top\}$，点$\bm y=(0,1)^\top$。

计算：$U$中的点写成$(t,0)^\top$，
\[
 \lVert\bm y-(t,0)^\top\rVert_\infty=\max\{|t|,1\}\geq1,
\]
等号当且仅当$|t|\leq1$时成立。

结论：$d(\bm y,U)=1$，所有满足$|t|\leq1$的点$(t,0)^\top$都是最近点，最近点不唯一。
\end{example}
内积诱导的范数下最近点唯一，并由正规方程算出（第~\ref{sec:ls-subspace-nearest-points}节）。

\subsection{完备性与Banach空间}
\label{sec:ls-banach-spaces}

收敛的定义要求事先知道极限$\bm x$。迭代算法、级数截断等过程却常常只给出一列越来越接近彼此的近似，极限是什么、在不在空间里，事先并不知道。第~\ref{chap:function-limits}章对实数列用柯西条件回答了这个问题；本小节把它搬到赋范空间，并说明它在有限维中总是可用。

\paragraph{柯西列与完备性}
赋范空间中的序列$\bm x_n$若满足：对任意$\varepsilon>0$，存在$N$使$m,n>N$时$\lVert\bm x_m-\bm x_n\rVert<\varepsilon$，就称为柯西列。收敛序列都是柯西列：若$\bm x_n\to\bm x$，则$\lVert\bm x_m-\bm x_n\rVert\leq\lVert\bm x_m-\bm x\rVert+\lVert\bm x-\bm x_n\rVert$。反过来，柯西列只检查序列内部，不引用任何候选极限；它是否收敛，取决于空间里有没有那个极限。若赋范空间中每个柯西列都收敛到空间中的某个元素，就称该空间\term{完备}（complete）；完备的赋范空间称为\term{Banach空间}（Banach space，也译作巴拿赫空间）。实数在通常距离下完备，这正是第~\ref{chap:function-limits}章的柯西收敛准则。

\begin{theorem}[有限维赋范空间的完备性]
\label{thm:ls-finite-dimensional-complete}
有限维赋范空间在任意范数下都完备，因而是Banach空间；它的每个子空间也完备。
\end{theorem}
\begin{proof}
设$\bm x^{(m)}$是柯西列。由范数等价，$\lVert\bm x^{(m)}-\bm x^{(n)}\rVert_c\leq\lVert\bm x^{(m)}-\bm x^{(n)}\rVert/m_0$，其中$m_0>0$是式~\eqref{eq:linear-algebra-finite-norm-equivalence}中的下界常数，所以每个坐标都是实柯西数列，由实数的完备性各有极限$x_j$。令$\bm x=\sum_jx_j\bm b_j$，有限个坐标的差同时趋于零，所以$\lVert\bm x^{(m)}-\bm x\rVert_c\to0$，再由范数等价得$\bm x^{(m)}\to\bm x$。子空间本身也是有限维赋范空间，同样完备；这也与定理~\ref{thm:ls-subspace-closed}一致：子空间中的柯西列在大空间中有极限，而极限留在闭的子空间里。
\end{proof}
证明的关键是有限维：坐标只有有限个，“每个坐标都收敛”就推出“最大坐标差趋于零”。坐标有无穷多个时这一步会失败：无限维赋范空间未必完备，是否完备取决于空间与范数的搭配。一般度量空间的完备性见\BookChapterRef{chap:metric-spaces}{第~\ref{chap:metric-spaces}章}，函数空间的完备性见\BookChapterRef{chap:functional-analysis-and-approximation}{第~\ref{chap:functional-analysis-and-approximation}章}。

完备性的用处在于，只要能证明近似序列是柯西列，就能断定极限存在并且仍在空间中，而不必事先找出极限。有限维中这一点由定理~\ref{thm:ls-finite-dimensional-complete}自动保证，所以第~\ref{chap:linear-analysis}章与第~\ref{chap:linear-equations-and-optimization}章讨论矩阵幂、迭代与级数的收敛时，只需估计误差，不必担心极限落到空间之外。范数等价只给出夹逼常数，用于在不同范数之间互换收敛结论，并不提供衰减速度：同一个收敛序列在一种范数下单调下降，在另一种范数下可能先上升，速度要在选定的范数中由递推结构给出。
